\begin{center}\hypertarget{native-228-article165}{}\pdfbookmark[2]{165.}{article-165}165.\end{center}
Das folgende Verfahren scheint das einfachste zu sein, um eine bereits annähernd bekannte Bahn an vier Beobachtungen anzupassen. Es seien $x$, $y$ die ungefähren Abstände des \SourceLeaf{229}Himmelskörpers von der Erde bei den beiden vollständigen Beobachtungen. Aus diesen Abständen berechnet man die entsprechenden heliozentrischen Orte und daraus die Elemente selbst. Aus diesen Elementen berechnet man dann die geozentrischen Längen oder Rektaszensionen für die beiden übrigen Beobachtungen. Stimmen diese mit den beobachteten Werten überein, so brauchen die Elemente nicht weiter korrigiert zu werden. Andernfalls notiert man die Differenzen $X$, $Y$ und wiederholt dieselbe Rechnung noch zweimal mit geringfügig veränderten Werten von $x$, $y$. So erhält man drei Sätze von Werten der Größen $x$, $y$ und der Differenzen $X$, $Y$. Nach den Anweisungen in Artikel 120 gewinnt man daraus die korrigierten Werte von $x$, $y$, denen die Werte $X=0$, $Y=0$ entsprechen. Eine entsprechende Rechnung auf der Grundlage dieses vierten Satzes liefert die Elemente, mit denen alle vier Beobachtungen richtig wiedergegeben werden.

Hat man die Wahl, so behält man vorzugsweise diejenigen Beobachtungen vollständig bei, aus denen sich die Lage der Bahn am genauesten bestimmen lässt. Das sind die erste und die letzte Beobachtung, sofern sie eine heliozentrische Bewegung von höchstens $90$ Grad umfassen. Sind die Beobachtungen jedoch nicht gleich genau, so lässt man die Breiten oder Deklinationen derjenigen beiseite, die man für weniger genau hält.
